\section{Корпус и его характеристики}

Тематика корпуса --- \textbf{кино и сериалы}: фильмы, сериалы, мультфильмы и
кинематографические персоналии. Источника два:

\begin{itemize}
    \item \href{https://ru.wikipedia.org/}{\textbf{ru.wikipedia.org}} --- энциклопедические
          статьи о фильмах, сериалах и кинематографистах;
    \item \href{https://www.kino-teatr.ru/}{\textbf{www.kino-teatr.ru}} --- карточки
          фильмов, биографии актёров, авторские статьи и обзоры.
\end{itemize}

\subsection*{Из чего состоит текст}

\textbf{Википедия}, кодировка UTF-8. Статьи собраны по устойчивому шаблону, разделы почти
всегда одни и те же: \emph{Сюжет}, \emph{В ролях}, \emph{Съёмочная группа},
\emph{Создание}, \emph{Критика}, \emph{Награды}, \emph{Примечания}, \emph{Литература},
\emph{Ссылки}. Три последних раздела --- списки источников и внешних адресов, к
содержанию документа они отношения не имеют.

\textbf{Кино-Театр.Ру}, кодировка \textbf{windows-1251}. Это первое, обо что
спотыкаешься: при чтении в UTF-8 весь текст превращается в мусор. Шаблона у статей нет,
это обычный авторский текст с подзаголовками; биографии и карточки фильмов устроены
проще --- несколько абзацев плюс фильмография списком.

\subsection*{Разметка}

Обе площадки отдают обычный HTML, текст виден прямо в исходном коде страницы. Основной
текст лежит в одном контейнере: у Википедии это
\texttt{<div class="mw-content-ltr mw-parser-output">}, у Кино-Театр.Ру ---
\texttt{<div class="art article\_text">}. Всё, что вне контейнера, --- меню, боковые
панели и подвал.

Внутри текста у Википедии встречаются карточка фильма
(\texttt{<table class="infobox">}) со структурированными полями, оглавление, сноски
\texttt{<sup class="reference">} и заголовки разделов \texttt{<h2>}.

\subsection*{Дополнительная мета-информация}

У Википедии: канонический адрес \texttt{<link rel="canonical">} и блок JS-конфигурации с
внутренним номером статьи \texttt{wgArticleId} и полным списком её категорий
\texttt{wgCategories} --- фактически готовая рубрикация. Даты документа в HTML нет.

У Кино-Театр.Ру мета-информации больше: набор Open Graph (\texttt{og:title},
\texttt{og:description}, \texttt{og:url}, \texttt{og:image}), микроразметка schema.org
через \texttt{itemprop} и точная дата публикации \texttt{datePublished}.

\pagebreak
